% =====================================================================
%  DRIFTGUARD — Section III: System Design.
%  Narrative rewrite (NOT a clean-up of reboot/docs/System_Design.tex).
%  Facts/numbers sourced from the bible; prose written fresh, story-first.
%  Method only; all measurements live in Section IV (\ref{sec:eval}) and
%  the safety layer in Section V (\ref{sec:safety}).
%
%  FIGURE: needs an architecture diagram (figures/architecture.pdf).
%  Placeholder included below — produce the diagram before submission.
% =====================================================================

\section{System Design}
\label{sec:design}

A phishing detector that must survive drift faces a tension. It has to answer
in milliseconds at the edge, yet the knowledge it answers with goes stale in
weeks. DRIFTGUARD resolves this by refusing to do everything in one place. A
small, fast model gives the verdict; a separate, label-free signal decides
\emph{when} the model has falle n behind; and only then, on a tight budget, a
large language model is consulted to manufacture the few gold labels needed to
bring the model back. The design mirrors how a security operations centre
actually runs --- cheap automated triage in line, expensive expert judgement
on demand --- and turns the LLM from an always-on dependency into a periodic
recalibration tool (Figure~\ref{fig:arch}).

Two principles run through the whole system and are worth stating up front.
First, \emph{honesty before recovery}: before any adaptation, we sterilize the
feature space of signals that look predictive only because of how the data was
collected, so that the accuracy we later recover is real and not an artefact.
Second, \emph{nothing informative is wasted}: the signals we forbid the model
from using are not discarded but handed to the oracle, which can exploit
contextually what a statistical model can only exploit as leakage. The rest of
this section develops the three components in turn; the operational safety
layer that wraps them is deferred to Section~\ref{sec:safety}.

\begin{figure}[t]
  \centering
  % TODO: produce figures/architecture.pdf (edge model -> pool -> drift
  % detector -> selection -> LLM oracle -> full-union retrain -> snapshot,
  % with the safety layer auditing the oracle stream).
  \includegraphics[width=\columnwidth]{figures/architecture.pdf}
  \caption{DRIFTGUARD as a closed loop. The leakage-clean detector serves
  instant verdicts at the edge and quietly buffers the samples it is unsure
  about. A label-free monitor watches the \emph{input} distribution and, when
  it has shifted enough, triggers a round: a small batch of uncertain,
  drift-aligned samples is sent to the LLM oracle for gold labels, and the
  detector is retrained on the full union of its seed corpus and every label
  acquired so far. A safety layer (Section~\ref{sec:safety}) audits the oracle
  stream throughout.}
  \label{fig:arch}
\end{figure}

\subsection{A Leakage-Clean Detection Core}
\label{ssec:core}

The detector is a light meta-learner that sits on top of two existing,
pre-trained phishing classifiers used as frozen \emph{proposers}: a graph
neural network over the rendered DOM and a gradient-boosted model over HTML
features. Neither is retrained inside the loop; each contributes a phishing
probability and the two are combined with simple agreement features. On their
own, the proposers are brittle --- under temporal drift the DOM-structure
signal essentially collapses to chance --- which is precisely why we do not
trust them blindly and instead ground the verdict in evidence they do not see.

That evidence comes from a \emph{passive OSINT measurement layer}. For every
page we measure the lifecycle of its domain from free, key-less public sources:
its popularity (the Tranco list), its archival history (the Wayback Machine),
its certificates (Certificate Transparency logs), its registration record
(RDAP), and its DNS and network footprint. We deliberately restrict ourselves
to sources that carry no label leakage by construction --- popularity,
archival, registration, routing --- and exclude anything correlated with the
phishing ground truth, such as blocklists, from the model's feature space. The
result is a few dozen factual lifecycle measurements per page, alongside the
proposer opinions and a handful of URL and HTML features.

The methodological centrepiece is what happens next. Many of these features are
strongly predictive on the training corpus yet betray the model under drift,
and the reason is subtle: a feature like ``number of Wayback snapshots''
separates the classes almost perfectly on historical data not because it
distinguishes phishing from benign, but because the historical phishing domains
are now \emph{dead} and therefore unarchived --- it measures ``alive versus
dead'', a temporal accident, not the phenomenon we care about. We catch such
features with a single statistical rule: any feature whose single-feature
discriminative power is suspiciously high on the training set \emph{and}
inconsistent across datasets is rejected as leakage. The same rule, applied
uniformly, turns out to catch two failures of completely different nature --
the temporally leaky OSINT signals \emph{and} the proposers' own
near-perfect-but-non-transferable probabilities --- under one criterion. We
quantify exactly which features it removes, and what that does to
cross-distribution accuracy, in Section~\ref{ssec:eval-gate}.

Crucially, a rejected feature is removed from the \emph{model}, not from the
\emph{system}: the gated-out lifecycle signals are preserved and forwarded to
the oracle, where contextual reasoning can use them safely. On the sterilized
feature space we deploy a logistic-regression detector. It starts from a
high-precision, low-recall operating point --- the pages it calls phishing are
almost always phishing, but under drift it misses many --- and those misses are
exactly the region the self-repair loop is built to recover.

\subsection{An LLM Oracle for Gold Labels}
\label{ssec:oracle}

When the loop needs a label it cannot get for free, it asks a large language
model. The oracle's job is narrow but demanding: given everything we know about
a single page, return a verdict reliable enough to retrain on. We give it three
complementary views of the page. The first is the cleaned HTML --- stripped of
styling and binary blobs but keeping the forensic markers (hidden fields,
password inputs, form targets, inline scripts) a verdict actually turns on. The
second is the OSINT dossier, including the very lifecycle signals the leakage
gate removed from the model: presented as plain facts (``domain registered 207
days ago; never archived; certificate rotated every 90 days''), they let the
oracle reason about throwaway phishing infrastructure that the statistical
model was rightly forbidden from exploiting. The third, when useful, is the
proposers' opinions, framed as advice rather than truth.

We ask the oracle to reason in fixed steps and to commit to a decision rather
than hedge, and we require it to list both the indicators it found and the
checks that came up empty. This last requirement is not cosmetic: the two
parallel lists are what later lets a verifier check each claim against the page,
turning an opaque verdict into an auditable one (Section~\ref{sec:safety}). The
output is a single structured record, robust to malformed generations, with the
verdict drawn from a closed taxonomy so that downstream analysis can reason
about \emph{why} a page was flagged, not just whether. Which combination of
evidence makes the best oracle is itself an empirical question, settled in
Section~\ref{ssec:eval-oracle}.

\subsection{The Self-Repair Loop}
\label{ssec:loop}

The loop ties the detector and the oracle together and is where the
``self-repair'' claim is earned. It must answer three questions without ever
seeing a ground-truth label at runtime: when has the model fallen behind, which
handful of pages are worth a costly oracle call, and how to fold the answers
back in without forgetting what the model already knew.

\paragraph{When: watching the input, not the confidence.}
The instinctive drift signal --- the model losing confidence on new data --- is
the wrong one here, and for an instructive reason. A leakage-clean,
well-calibrated detector does \emph{not} get hesitant under drift; it stays
confident and is simply wrong, making confident mistakes on inputs it has never
seen. Any trigger that waits for confidence to collapse therefore waits
forever. DRIFTGUARD instead measures the shift in the \emph{input}
distribution: it compares the running distribution of the model's predictions
against a reference taken at deployment, and fires when that divergence grows.
We keep this timing mechanism deliberately simple --- in most runs, a scheduled
batching policy --- precisely so that the gains we report can be attributed to
\emph{which} samples we label rather than to a bespoke trigger. The operational
consequences of this choice, including the cost of reacting late, are measured
in Section~\ref{ssec:eval-delay}.

\paragraph{Which: spending the budget where the drift is.}
When a round fires, the detector quietly maintains a pool of the samples it was
uncertain about, and we must pick a small batch to label. Classical active
learning picks the most uncertain points; but ``uncertain'' is defined relative
to the current model, blind to \emph{where the data has actually moved}. Our
selector, \emph{drift-anchored} sampling, corrects this. It scores a candidate
by its boundary uncertainty, then up-weights it by how far it has drifted from
the training distribution along the dimensions that have drifted most:
\[
s(x) \;=\; \underbrace{\frac{1}{|p(x)-0.5|+\epsilon}}_{\text{uncertainty}}
\;\cdot\;
\Big(1 + \alpha \textstyle\sum_i \delta_i\,|x_i-\mu_i|\Big),
\]
where $\delta_i$ is the measured per-feature drift and $\mu_i$ the training
mean. The effect is to spend labels on pages that are both hard \emph{and}
explanatory of the observed shift, degrading gracefully to plain uncertainty
sampling when nothing has drifted. We benchmark it against the standard active-
learning strategies, and surface where it helps and where it trades off, in
Section~\ref{ssec:eval-recovery}; the single weight $\alpha$ is shown not to
matter in Section~\ref{ssec:eval-alpha}.

\paragraph{How: retraining without forgetting.}
Each round, the detector is refit from scratch on the \emph{full union} of its
original corpus and every label acquired so far. This is the blunt but reliable
guarantee against catastrophic forgetting: because all of the original evidence
is re-presented at every fit, recovery on the drifted distribution cannot come
at the price of the original one --- a claim we verify rather than assume
(Section~\ref{ssec:eval-recovery}). The feature scaler, fit once at deployment,
is frozen for the life of the run, since re-estimating it on a few dozen new
samples each round would induce far more coefficient drift than it removes.
After each retrain the uncertainty pool is re-scored under the new model, so
pages that have become easy graduate out and the pool stays aligned with the
model's \emph{current} blind spots.

\medskip\noindent
This is the whole loop: an honest baseline, a contextual oracle fed the very
signals the baseline could not safely use, and a budgeted, drift-aware
acquisition cycle that retrains without forgetting. What we have not yet
addressed is how to run it safely when the oracle is an LLM and no ground truth
is available to check it --- the subject of Section~\ref{sec:safety}.
